\documentclass{article}
\usepackage{graphicx} % Required for inserting images
\usepackage{amsmath}
\usepackage{setspace}
\onehalfspacing
\usepackage{tikz}
\usepackage{pgfplots}
\pgfplotsset{compat=1.18}

\title{CS556-HW3 - Vector Spaces and Subspaces}
\author{John Rizzo}
\date{October 5, 2025}

\usepackage[margin=2cm]{geometry}

\begin{document}

\maketitle
\noindent
\textbf{Question 1} (10 points) Write the complete solution of the following linear system:

$x + 2y - z = 1$

$3x + 5y + 2z = 3$

$2x + y + 13z = 2$

Solution 

$\begin{bmatrix}
1 & 2 & -1 \\
3 & 5 & 2\\
2 & 1 & 13
\end{bmatrix}\begin{bmatrix}
x \\ y \\ z    
\end{bmatrix} = \begin{bmatrix}
    1 \\ 3 \\ 2
\end{bmatrix}$

R2 = R2 - 3R1

$\begin{bmatrix}
1 & 2 & -1 \\
0 & -1 & 5\\
2 & 1 & 13
\end{bmatrix}\begin{bmatrix}
x \\ y \\ z    
\end{bmatrix} = \begin{bmatrix}
    1 \\ 0 \\ 2
\end{bmatrix}$ 

 R3 = R3 - 2R1

$\begin{bmatrix}
1 & 2 & -1 \\
0 & -1 & 5\\
0 & -3 & 15
\end{bmatrix}\begin{bmatrix}
x \\ y \\ z    
\end{bmatrix} = \begin{bmatrix}
    1 \\ 0 \\ 0
\end{bmatrix}$

R2 = -1 R2

$\begin{bmatrix}
1 & 2 & -1 \\
0 & 1 & -5\\
0 & 0 & 0
\end{bmatrix}\begin{bmatrix}
x \\ y \\ z    
\end{bmatrix} = \begin{bmatrix}
    1 \\ 0 \\ 0
\end{bmatrix}$ RREF

Infinite solutions since the last row is all 0's.

\vspace{2em}
\noindent
\textbf{Question 2} (10 points) Find the rank of the following matrix:

$\begin{bmatrix}
0 & 1 & 2 & 1 \\
1 & 2 & 3 & 2 \\
3 & 1 & 1 & 3 
\end{bmatrix}$ 

Solution

Swap R1, R2

$\begin{bmatrix}
    1 & 2 & 3 & 2 \\
    0 & 1 & 2 & 1 \\
    3 & 1 & 1 & 3
\end{bmatrix}$

R3 = R3 - 3R1

$\begin{bmatrix}
    1 & 2 & 3 & 2 \\
    0 & 1 & 2 & 1 \\
    0 & -5 & -8 & 3
\end{bmatrix}$ 

 R3 = R3 + 5R2

$\begin{bmatrix}
    1 & 2 & 3 & 2 \\
    0 & 1 & 2 & 1 \\
    0 & 0 & 2 & 2 
\end{bmatrix}$

R2 = $1/2R2$

$\begin{bmatrix}
    \textbf{1} & 2 & 3 & 2 \\
    0 & \textbf{1} & 2 & 1 \\
    0 & 0 & \textbf{1} & 1
\end{bmatrix}$

\textbf{Rank is 3}

\vspace{2em}
\noindent
\textbf{Question 3} (10 points) Construct a matrix $A$ whose column space contains vectors $\begin{bmatrix}
    3 \\ 6 \\ 2
\end{bmatrix}$ and $\begin{bmatrix}4 \\ 0 \\ 1\end{bmatrix}$, and whose null space contains the vector $\begin{bmatrix}
    2 \\ 2 \\ 1
\end{bmatrix}$.

Solution

$A = \begin{bmatrix}
    3 & 4 & a \\
    6 & 0 & b \\
    2 & 1 & c
\end{bmatrix}$

$x = \begin{bmatrix}
    2 \\ 2 \\ 1
\end{bmatrix}$

Null Space, $Ax=0$

6 + 8 + 1a = 0, a = -14

12 + 0 + 1b = 0, b = -12

4 + 2 + c = 0, c = -6

So $A = \begin{bmatrix}
    3 & 4 & -14 \\
    6 & 0 & -12 \\
    2 & 1 & -6
\end{bmatrix}$ is a matrix with the two column vectors and a Null Space which contains the vector $\begin{bmatrix}
    2 \\ 2 \\ 1
\end{bmatrix}$

\vspace{2em}
\noindent
\textbf{Question 4} (10 points) Compute the following matrix-vector multiplication as:

$\begin{bmatrix}
    2 & 1 & 3 \\
    7 & 1 & 0 \\
    3 & 5 & 9
\end{bmatrix}\begin{bmatrix}
    2 \\ 4 \\ 1
\end{bmatrix}$

a) Linear combination of columns

$2 \times \begin{bmatrix}
    2 \\ 7 \\ 3
\end{bmatrix} + 4 \times \begin{bmatrix}
    1 \\ 1 \\ 5
\end{bmatrix} + 1 \times \begin{bmatrix}
    3 \\ 0 \\ 9
\end{bmatrix} = \begin{bmatrix}
    11 \\ 18 \\ 35
\end{bmatrix}$

b) Dot product of rows

$\begin{bmatrix}
    2 & 1 & 3 \\
    7 & 1 & 0 \\
    3 & 5 & 9
\end{bmatrix}\begin{bmatrix}
    2 \\ 4 \\ 1
\end{bmatrix} = \begin{bmatrix}
    11 \\ 18 \\ 35
\end{bmatrix}$

\vspace{2em}
\noindent
\textbf{Question 5} (10 points) Find the value of $k$ for which the matrix has:

$A = \begin{bmatrix}
    1 & 3 & 2 \\
    2 & 3 & 1 \\
    4 & 8 & k
\end{bmatrix}$

a) Dependent Columns

det(A) = $1 \begin{vmatrix}
    3 & 1 \\ 8 & k
\end{vmatrix} - 3 \begin{vmatrix}
    2 & 1 \\ 4 & k
\end{vmatrix} + 2 \begin{vmatrix}
    2 & 3 \\ 4 & 8
\end{vmatrix}$

            $= (3k - 8) - 3(2k - 4) + 2(4)$
            
            $ = 3k - 8 -6k + 12 + 8$
            
            $= -3k + 12$
            
            $k = \frac{-12}{-3} = 4$

For the columns to be dependent the determinate should be 0.  Setting k = 4 causes the determinate to be 0.

b) Independent Columns

Choosing any number other than 4 would cause the determinate not to be 0 and therefore the columns would be independent.

\vspace{2em}
\noindent
\textbf{Question 6} (20 points) Find a basis for the four fundamental subspaces of the matrix.

$A = \begin{bmatrix}
    0 & 1 & 2 & 3 & 4 \\
    0 & 1 & 2 & 4 & 6 \\
    0 & 0 & 0 & 1 & 2
\end{bmatrix}$

Null Space $N(A)$, ($Ax=0$)

$\begin{bmatrix}
    0 & 1 & 2 & 3 & 4 \\
    0 & 1 & 2 & 4 & 6 \\
    0 & 0 & 0 & 1 & 2
\end{bmatrix}x = \begin{bmatrix}
    0 \\ 0 \\ 0
\end{bmatrix}$

R2 = R2 - R1

$\begin{bmatrix}
    0 & 1 & 2 & 3 & 4 \\
    0 & 0 & 0 & 1 & 2 \\
    0 & 0 & 0 & 1 & 2
\end{bmatrix}x = \begin{bmatrix}
    0 \\ 0 \\ 0
\end{bmatrix}$

R3 = R3 - R2

$\begin{bmatrix}
    0 & 1 & 2 & 3 & 4 \\
    0 & 0 & 0 & 1 & 2 \\
    0 & 0 & 0 & 0 & 0 \\
    F & P & F & P & F
\end{bmatrix}x = \begin{bmatrix}
    0 \\ 0 \\ 0
\end{bmatrix}$ 

Column Space $C(A)$, ($Ax = b$)

The basis for the column space is $\begin{bmatrix}
    1 \\ 1 \\ 0
\end{bmatrix}$ and $\begin{bmatrix}
    4 \\ 6 \\ 2
\end{bmatrix}$.

Row Space = $R(A) = C^T$

The non-zero rows form the row space, $\begin{bmatrix}
    0 & 1 & 2 & 3 & 4
\end{bmatrix}$ and $\begin{bmatrix}
    0 & 1 & 2 & 4 & 6
\end{bmatrix}$

Null Space = N(A), ($Ax = 0$)

Let $\hat{x} = \begin{bmatrix}
    x_1, x_2, x_3, x_4, x_5
\end{bmatrix}^T$

From Row 2: $x_4 + 2x_5 = 0$ and so $x_4 = -2x_5$

From Row 1: $x_2 + 2x_3 + 3x_4 + 4x_5 = 0$, Substituting $x_4$ into the equation yields 

$x_2 + 2x_3 + 3(-2x_5) + 4x_5 = 0$

$x_2 + 2x_3 - 6x_5 + 4x_5 = 0$

$x_2 + 2x_3 - 2x_5 = 0$

$x_2 = -2x_3 + 2x_5$

Let $x_1 = t, x_3 = s, x_5 = r$

$\hat{x} = \begin{bmatrix}
    t \\ -2s + 2r \\ s \\ -2r \\ r
\end{bmatrix} = t\begin{bmatrix}
    1 \\ 0 \\ 0 \\ 0 \\ 0
\end{bmatrix} + s\begin{bmatrix}
    0 \\ -2 \\ 1 \\ 0 \\ 0
\end{bmatrix} + r\begin{bmatrix}
    0 \\ 2 \\ 0 \\ -2 \\ 1
\end{bmatrix}$

so the basis for the null space is $\begin{bmatrix}
    1 \\ 0 \\ 0 \\ 0 \\ 0
\end{bmatrix}, \begin{bmatrix}
    0 \\ -2 \\ 1 \\ 0 \\ 0
\end{bmatrix}, \begin{bmatrix}
    0 \\ 2 \\ 0 \\ -2 \\ 1
\end{bmatrix}$

Left Null Space = $N(A)^T$

$A^T = \begin{bmatrix}
    0 & 0 & 0 \\
    1 & 1 & 0 \\
    2 & 2 & 0 \\
    3 & 4 & 1 \\
    4 & 6 & 2
\end{bmatrix}$

Let $\hat{y} = \begin{bmatrix}
    y_1, y_2, y_3
\end{bmatrix}^T$

$0y_1 + 0y_2 + 0y_3 = 0$

$y_1 + y_2 = 0$, or $y_1 = -y_2$

$2y_1 + 2y_2 + 0y_3 = 0$ or $y_1 = -y_2$

$3y_1 + 4y_2 + 1y_3 = 0$ or when substituting $-y_2$ for $y_1$ $-3y_2 + 4y_2 + y_3 = y_2 + y_3 = 0$ or $y_3 = -y_2 = y_1$

$4y_1 + 6y_2 + 2y_3 = 0$

Finally, if we let $y_2 = s$, $y_1 = -s$, $y_3 = -s$ we have the basis of $\begin{bmatrix}
    -1 \\ 1 \\ -1
\end{bmatrix}$

\vspace{2em}
\noindent
\textbf{Question 7} (10 points) Consider the vectors $\hat{v} = \begin{bmatrix}
    1 \\ 2
\end{bmatrix}$ and $\hat{w} = \begin{bmatrix}
    1 \\ 0
\end{bmatrix}$
\noindent
a) In the x-y plane mark all nine linear combinations $c\hat{v} + d\hat{w}$, with $c = \{-2, 0 -2\}$ and $d = \{ 0, 1, 2\}$.

\noindent
$-2\hat{v} + 0\hat{w}$  (-2, 0) \\
$-2\hat{v} + 1\hat{w}$  (-2, 1) \\
$-2\hat{v} + 2\hat{w}$  (-2, 2) \\
$0\hat{v} + 0\hat{w}$   (0, 0) \\
$0\hat{v} + 1\hat{w}$   (0, 1) \\
$0\hat{v} + 2\hat{w}$   (0, 2) \\
$2\hat{v} + 0\hat{w}$   (2, 0) \\
$2\hat{v} + 1\hat{w}$   (2, 1) \\
$2\hat{v} + 2\hat{w}$   (2, 2)\\

$\begin{tikzpicture}
  \begin{axis}[
    xlabel={$x$-axis}, ylabel={$y$-axis},
    width=10cm, height=10cm,
    grid=major,
  ]
  \addplot[
    only marks,
    mark=*,
    color=blue]
  coordinates {
    ((-2,0) (-2,1)(-2,2)(0,0)(0,1)(0,2)(2,0)(2,1)(2,2))
  };
  \end{axis}
\end{tikzpicture}$

\noindent
b) What shape do all linear combinations $c\hat{v} + d\hat{w}$ fill? A line? The whole plane? Are the vectors $\hat{v}$ and $\hat{w}$ independent?

The linear combinations span the whole plane.  

From computing the determinate will can assess the independence of the vectors $\hat{v}$ and $\hat{w}$.

$det\begin{bmatrix}
    1 & 1 \\
    2 & 0
\end{bmatrix} = (1 \times 0 - 1 \times 2) = -2$

The vectors are independent because the determinate is not 0.

\vspace{2em}
\noindent
\textbf{Question 8} (10 points) Consider the vectors $\hat{u} = \begin{bmatrix}
    1 \\ 0 \\ 1
\end{bmatrix}$, $\hat{v} = \begin{bmatrix}
    0 \\ -1 \\ 1
\end{bmatrix}$ and $\hat{w} = \begin{bmatrix}
    2 \\ 1 \\ 0
\end{bmatrix}$.

a) Can you solve the system $x\hat{u} + y\hat{v} + z\hat{w} = \hat{b}$ if $\hat{b} = \begin{bmatrix}
    0 \\ 0 \\ 1
\end{bmatrix}$?

$\begin{bmatrix}
    1 & 0 & 2 & 0 \\
    0 & -1 & 1 & 0 \\
    1 & 1 & 0 & 1
\end{bmatrix}$

R3 = R3 - R1

$\begin{bmatrix}
    1 & 0 & 2 & 0 \\
    0 & -1 & 1 & 0 \\
    0 & 1 & -2 & 1
\end{bmatrix}$

R3 = R3 + R2

$\begin{bmatrix}
    1 & 0 & 2 & 0 \\
    0 & -1 & 1 & 0 \\
    0 & 0 & -1 & 1
\end{bmatrix}$

-1R2, -1R3

$\begin{bmatrix}
    1 & 0 & 2 & 0 \\
    0 & 1 & -1 & 0 \\
    0 & 0 & 1 & -1
\end{bmatrix}$

R2 = R2 + R3

$\begin{bmatrix}
    1 & 0 & 2 & 0 \\
    0 & 1 & 0 & -1 \\
    0 & 0 & 1 & -1
\end{bmatrix}$

R1 = R1 -2R3

$\begin{bmatrix}
    1 & 0 & 0 & 2 \\
    0 & 1 & 0 & -1 \\
    0 & 0 & 1 & -1
\end{bmatrix}$

$x = 2$
$y = -1$
$z = -1$

b) What if $\hat{b} = \begin{bmatrix}
    0 \\ 0 \\ 0
\end{bmatrix}$? How many solutions are there?

$\begin{bmatrix}
    1 & 0 & 2 & 0 \\
    0 & -1 & 1 & 0 \\
    1 & 1 & 0 & 0
\end{bmatrix}$

R3 = R3 - R1

$\begin{bmatrix}
    1 & 0 & 2 & 0 \\
    0 & -1 & 1 & 0 \\
    0 & 1 & -2 & 0
\end{bmatrix}$

R3 = R3+R2

$\begin{bmatrix}
    1 & 0 & 2 & 0 \\
    0 & -1 & 1 & 0 \\
    0 & 0 & -1 & 0
\end{bmatrix}$

-1R2, -1R2

$\begin{bmatrix}
    1 & 0 & 2 & 0 \\
    0 & 1 & -1 & 0 \\
    0 & 0 & 1 & 0
\end{bmatrix}$

R2 = R2+R3

$\begin{bmatrix}
    1 & 0 & 2 & 0 \\
    0 & 1 & 0 & 0 \\
    0 & 0 & 1 & 0
\end{bmatrix}$

R1 = R1-2R3

$\begin{bmatrix}
    1 & 0 & 0 & 0 \\
    0 & 1 & 0 & 0 \\
    0 & 0 & 1 & 0
\end{bmatrix}$

\textbf{\textit{There are infinitely many solutions?}}

c) Are the vectors $\hat{u}$, $\hat{v}$, and $\hat{w}$ dependent or independent?

d) Use parts (a) - (c) to decide if $A = \begin{bmatrix}
    1 & 0 & 2 \\
    0 & -1 & 1 \\
    1 & 1 & 0
\end{bmatrix}$ is an invertible matrix or not.

\vspace{2em}
\noindent
\textbf{Question 9} (10 points) Consider the linear system for some constants $b$ and $g$:

$x - 2y + 3z = 3$

$2x + y + bz = -4$

$x + 0y + 1z = g$

Converted to an augmented matrix

$\left[
\begin{array}{ccc|c}
  1 & -2 & 3 & 3 \\
  2 & 1 & b& -4\\
  1 & 0 & 1 & g
\end{array}
\right]$

Gaussian Elimination

R2 = R2 - 2R1

$\left[
\begin{array}{ccc|c}
  1 & -2 & 3 & 3 \\
  0 & 5 & b-6 & -10\\
  1 & 0 & 1 & g
\end{array}
\right]$

R3 = R3 - R1

$\left[
\begin{array}{ccc|c}
  1 & -2 & 3 & 3 \\
  0 & 5 & b-6 & -10\\
  0 & 2 & -2 & g - 3
\end{array}
\right]$

$R3 = R3 - \frac{2}{5}R2$

$\left[
\begin{array}{ccc|c}
  1 & -2 & 3 & 3 \\
  0 & 5 & b-6 & -10\\
  0 & 0 & \frac{2}{5} (1 - b) & g + 1
\end{array}
\right]$

a) What constant b makes the system singular (missing a pivot).

When the third row's coefficient for z is zero.  This is the case when $\frac{2}{5}(1-b) = 0$, which is the case when $1-b = 0$, $b=1$.

b) For the value of $b$ found in Part (a), for which values of $g$, does the system have infinitely many solutions?

Continuing for $b=1$ we can see the values of $g$ for which the system has infinitely many solutions.

$\left[
\begin{array}{ccc|c}
  1 & -2 & 3 & 3 \\
  0 & 5 & b-6 & -10\\
  0 & 0 & 0 & g + 1
\end{array}
\right]$

For there to be infinitely many solutions the third row solution should be 0 = 0 which is the case when g = -1.

c) Find two distinct solutions of the system for that $g$.

For the values $b=1$ and $g = -1$ the system simplifies to

$x - 2y + 3z = 3$

$2x + y + 1z = -4$

$x+z = -1$ or $x = -1 -z$

Substituting $x=-1-z$ back into the other two equations yields

Substitution 1

$(-1 - z) - 2y + 3z = 3$

$- 2y + 2z = 4$

$-y+z=2$

$z-y = 2$

Substitution 2

$2(-1 -z) + y + z = -4$

$y - z = -2$

$y = z - 2$

so we have

$x=-1-z$ 

$y = z - 2$

z is free

Solutions

For $z = 0$, $x=-1-0=-1$, $y = 0 - 2 = -2$

For $z = 1$, $x = -1 -1 = -2$, $y = 1 -2 = -1$


\end{document}
